% Meeting 23
	\pagebreak
	\section*{Mar. 17, 2026}
	\addcontentsline{toc}{section}{Meeting 24}

	\textbf{C.A.V.E Attendees:} Abdul, Andrew, Nicholas, Dr. Giamou\\
	\textbf{Notetaker:} Nicholas

	\subsection*{Meeting with Dr. Giamou}
		\subsubsection*{Agenda}
			\begin{itemize}
				\item Post Rev1 catch up and caving information
			\end{itemize}

		\subsubsection*{Minutes}
			\begin{itemize}
				\item getting forms together may not come in time for approval
				\item conditions may not be favourable for caving on march 23 given the cooling, will have a melt
				\item SHowing Giamou slides from the rev1 presentation to update on project status
				\item issues with recording, fusion
				\item TOFs have low resolution, poor connection stability (likely due to cost)
				\item Worried about being able to collect cave data
				\item Have GTSAM working with IMU data, can use either rotations or add a preintegration with velocity, but with kinda of poor results
				\item Manual loop closure still
				\item Filtering is iffy sometimes, removes excessive amounts of data
				\item Pyramidal effect around cone of each point capture
				\item Once we have cave data, may need to add more manual intervention to tell ICP what results are good or bad
				\item Giamou agrees on most of demo contract, 1.1 most likely to be difficult. 1.3 remains to be seen
				\item For a capstone demo, we can do a small room, show that cave worked or didnt work to a certain amount
				\item Experimental/exploratory, learning from the design and we can present these findings
				\item Action items:
				\item March 27 poster deadline
				\item can add stuff on top of it later, or show on the laptop during expo
				\item reminded that we cant go to caves alone (without Giamou)
				\item want to integrate everything on the helmet on thursday (mar 19)
				\item At this point, more human intervention is needed, or tuning uncertainties/improving preintegration
				\item Due to offline processing, can randomly initialize. Currently initializing poses to identity, could use open loop ICP estimate
				\item More sophisticated would be to form optztn problem ass least squares, drop rotation matrices, solve linear system
				\item sensor quality likely limiting factor at this point
				\item Covariance matrix for ICP currently random, theres a paper defining it but we could do better- datasheet may provide values for pointcloud
				\item sensitivity analysis method that adds noise and compares to original ICP result 
			\end{itemize}

		\subsubsection*{Action Items}
			\begin{itemize}
				\item Final integration and testing in a basic environment
				\item Finish documentation, start on poster
				\item Go caving
			\end{itemize}

